% SPDX-License-Identifier: CC-BY-4.0
\section{Discussion}
\label{sec:discussion}

\subsection{When \sediment{} helps cleanly}

\sediment{} is at its strongest when the benign distribution is
broadly Gaussian per feature, the adversary budget $\budget$ is at
most $\trim/2$, and the target attack behavior $\target$ sits
outside the typical benign envelope. Under those conditions
\cref{prop:robust} applies: the trimmed mean is provably resistant
up to the breakdown rate of $\trim/2$, and the Gaussian-truncation
correction restores an unbiased $\sigma$ estimate. The
clean-condition false-positive rate matches the naive estimator
(\cref{sec:eval-fp-cost}), which makes \trim{} a near-free knob
to dial up against adversaries. Choosing \trim{} larger than
necessary costs little. Choosing it smaller than $2\budget$ costs a
lot.

\subsection{When the bias correction over-trims}

The closed-form correction $c(\trim)$ assumes a unimodal Gaussian
null. On bimodal benign distributions the trimmed middle has
strictly lower variance than the full mixture, so the corrected
$\sigma$ underestimates the true population spread. The result,
visible in \cref{fig:exp06}, is inflated absolute scores at low
budgets even though threshold-invariant discriminability ratios
stay reasonable. The mitigation for production deployment is
workload-specific. Operators whose feature distributions are known
to be multimodal should pair \sediment{} with a mixture-aware
preprocessor or fall back to a different estimator family for
those features.

\subsection{Concept drift versus the maintenance window}

\sediment{} addresses adversarial drift, not concept drift. The
\texttt{maintenance start/end} CLI subcommands give operators a way
to mark intervals where baseline shifts are expected. Maintenance
windows are written to the audit log and can be queried during
incident review. The two mechanisms are orthogonal. \sediment{}
tightens what counts as a robust baseline at fit time. Maintenance
windows tell the operator-facing review path which alerts to expect
during legitimate transitions.

\subsection{Residual risk from adaptive adversaries}

The white-box experiment in \cref{sec:eval-white-box} shows an
adversary who knows \trim{} and grid-searches the optimal placement
does not gain a meaningful advantage over the at-target injection.
This is not a security proof. An adaptive adversary with control
over the full attack-fit fixed-point loop, an oracle for the
post-fit $(\mu, \sigma)$ moments, or the ability to shape multiple
correlated features in concert may yet find a placement that
outperforms target injection. We report the negative result of our
grid search as evidence that the obvious attack class does not
break the construction. We do not claim the construction is
unconditionally robust.
